\chapter{Stochastyczne równania różniczkowe ze skokami}
\label{ch:jump-sde}

\section{Motywacja: dlaczego zwykłe SDE nie opisuje nagłych zmian?}

Model $\dd X_t=b(X_t)\,\dd t+\sigma(X_t)\,\dd W_t$ ma ciągłe
trajektorie. Nie opisuje zatem nagłego uszkodzenia elementu, wpłaty
kapitału ani natychmiastowego przyrostu liczby cząstek. Dodanie
skokowego składnika prowadzi do procesu càdlàg. Jego wartości
bezpośrednio przed skokiem i bezpośrednio po skoku nie są tym samym:
$X_{t-}$ zawiera informację dostępną przed skokiem, a $X_t$ ---
już wynik skoku. Dlatego poprawne równanie ma po lewej stronie
$\dd X_t$, ale w prawych współczynnikach skokowych pojawia się
$X_{t-}$.

Są tu dwie naturalne konwencje. Całka względem $N$ dodaje
rzeczywiste skoki, natomiast całka względem
$\widetilde N=N-\dd t\,\dd\nu$ odejmuje od nich średni dryf.
Pozornie kosmetyczny wybór wpływa na współczynnik dryfu, lecz
nie zmienia skoków procesu.

\begin{definicja}[Równanie skokowe i jego silne rozwiązanie]
\label{def:jump-sde}
Niech $T>0$, $E$ będzie przestrzenią polską z miarą
$\sigma$-skończoną $\nu$, $E=E_0\sqcup E_1$ będzie podziałem
na zbiory borelowskie, gdzie $\nu(E_1)<\infty$.
Na przestrzeni z filtracją spełniającą warunki zwykłe niech $W$
będzie $m$-wymiarowym procesem Wienera, a $N$ losową miarą
Poissona o intensywności $\dd t\,\dd\nu(z)$, przy czym
$(W,N)$ ma niezależne przyrosty względem wspólnej filtracji.
Niech $x\in\R^d$ oraz
\[
\begin{gathered}
b:[0,T]\times\R^d\to\R^d,\quad
\sigma:[0,T]\times\R^d\to\R^{d\times m},\\
\gamma:[0,T]\times\R^d\times E_0\to\R^d,\quad
\Gamma:[0,T]\times\R^d\times E_1\to\R^d
\end{gathered}
\]
będą funkcjami borelowskimi. \emph{Silnym rozwiązaniem} równania
\[
\begin{split}
X_t=x&+\int_0^t b(s,X_{s-})\,\dd s
 +\int_0^t\sigma(s,X_{s-})\,\dd W_s\\
&+\int_{(0,t]\times E_0}\gamma(s,X_{s-},z)\,\widetilde N(\dd s,\dd z)
 +\int_{(0,t]\times E_1}\Gamma(s,X_{s-},z)\,N(\dd s,\dd z)
\end{split}
\tag{13.1}\label{eq:jump-sde}
\]
nazywamy proces adaptowany càdlàg $X$ o $X_0=x$, dla którego
wszystkie całki istnieją we wskazanym sensie, a tożsamość zachodzi
prawie na pewno jednocześnie dla wszystkich $t\in[0,T]$.
Przez $\|\sigma\|_{\mathrm{HS}}$ oznaczamy normę Hilberta--Schmidta.
\end{definicja}

\begin{uwaga}[Dwie skale skoków]
Na $E_1$ występuje skończenie wiele punktów miary $N$ na każdym
przedziale zwartym, zatem składnik z $\Gamma$ jest zwykłą losową
sumą. Zbiór $E_0$ może mieć nieskończoną miarę, a małe skoki
trzeba sumować po kompensacji. Dla samego istnienia procesu
skokowego nie żądamy drugiego momentu $\Gamma$; w pierwszym
twierdzeniu przyjmiemy go celowo, aby podać zupełny dowód
metodą Picarda w przestrzeni $L^2$.
\end{uwaga}

\section{Przygotowanie: zamiana dużych skoków na całkę skompensowaną}

\begin{lemat}[Przepisanie równania w klasie $L^2$]
\label{lem:rewrite-jump-sde}
Załóżmy, że $\nu(E_1)<\infty$ i dla każdego $(t,x)$ zachodzi
$\int_{E_1}\|\Gamma(t,x,z)\|^2\,\dd\nu(z)<\infty$.
Oznaczmy
\[
 \beta(t,x)=b(t,x)+\int_{E_1}\Gamma(t,x,z)\,\dd\nu(z),\qquad
 h(t,x,z)=\gamma(t,x,z)\ind{E_0}(z)+\Gamma(t,x,z)\ind{E_1}(z).
\]
Jeśli obie strony mają sens całkowy, równanie
\eqref{eq:jump-sde} jest równoważne równaniu
\[
 X_t=x+\int_0^t\beta(s,X_{s-})\,\dd s
 +\int_0^t\sigma(s,X_{s-})\,\dd W_s
 +\int_{(0,t]\times E}h(s,X_{s-},z)\,\widetilde N(\dd s,\dd z).
 \tag{13.2}\label{eq:jump-sde-comp}
\]
\end{lemat}
\begin{proof}
Nierówność Cauchy'ego--Schwarza daje
\[
\int_{E_1}\|\Gamma(t,x,z)\|\,\dd\nu(z)
\le \nu(E_1)^{1/2}
\left(\int_{E_1}\|\Gamma(t,x,z)\|^2\,\dd\nu(z)\right)^{1/2}.
\]
Można więc odjąć skończony kompensator dużych skoków:
\[
\int_{(0,t]\times E_1}\Gamma\,\dd N
=\int_{(0,t]\times E_1}\Gamma\,\dd\widetilde N
 +\int_0^t\int_{E_1}\Gamma\,\dd\nu\,\dd s.
\]
Tożsamość sprawdzamy najpierw dla procesów prostych,
a następnie dla ogólnych współczynników przez lokalizację
i izometrię z rozdziału~\ref{ch:poisson-random-measure}.
Zsumowanie składników dryfu daje \eqref{eq:jump-sde-comp}.
\end{proof}

\begin{przyklad}[Różnica pomiędzy dryfem a średnim przyrostem]
Niech $N$ będzie procesem Poissona o intensywności
$\lambda>0$. Dla $X_t=x+bt+cN_t$ mamy
$\E X_t=x+(b+\lambda c)t$. Jeśli zamiast tego zapiszemy
$X_t=x+\widetilde b t+c(N_t-\lambda t)$, to oba opisy
są identyczne dokładnie przy $\widetilde b=b+\lambda c$.
Nie należy więc porównywać współczynników przy $\dd t$ bez
sprawdzenia, czy skoki zostały skompensowane.
\end{przyklad}

\section{Istnienie, jednoznaczność i dokładny dowód metodą Picarda}

\begin{twierdzenie}[Silne istnienie i jednoznaczność w klasie $L^2$]
\label{th:jump-picard}
Niech dane będą jak w Definicji~\ref{def:jump-sde}.
Załóżmy, że istnieją $L,K_0<\infty$ takie, iż dla wszystkich
$t\in[0,T]$ i $x,y\in\R^d$ zachodzą
\[
\begin{split}
&\|b(t,x)-b(t,y)\|^2+
 \|\sigma(t,x)-\sigma(t,y)\|_{\mathrm{HS}}^2\\
&\quad+\int_{E_0}\|\gamma(t,x,z)-\gamma(t,y,z)\|^2\dd\nu(z)
 +\int_{E_1}\|\Gamma(t,x,z)-\Gamma(t,y,z)\|^2\dd\nu(z)
 \le L\|x-y\|^2,
\end{split}
\tag{13.3}\label{eq:jump-lip}
\]
oraz
\[
\sup_{t\in[0,T]}\left(
\|b(t,0)\|^2+\|\sigma(t,0)\|_{\mathrm{HS}}^2
+\int_{E_0}\|\gamma(t,0,z)\|^2\dd\nu(z)
+\int_{E_1}\|\Gamma(t,0,z)\|^2\dd\nu(z)
\right)\le K_0.
\tag{13.4}\label{eq:jump-at-zero}
\]
Wtedy istnieje adaptowane silne rozwiązanie càdlàg $X$
równania \eqref{eq:jump-sde}, jedyne z dokładnością
do nierozróżnialności, i
$\E\sup_{t\le T}\|X_t\|^2<\infty$.
Dla dwóch rozwiązań napędzanych tą samą parą $(W,N)$
i z tym samym warunkiem początkowym zachodzi
$\Prob(X_t=Y_t\ \text{dla wszystkich }t\le T)=1$.
\end{twierdzenie}
\begin{proof}
\emph{Krok 1: kontrola współczynników po kompensacji.}
Niech $\beta,h$ będą z Lematu~\ref{lem:rewrite-jump-sde}.
Z \eqref{eq:jump-lip} i nierówności Cauchy'ego--Schwarza
wynika, z pewną stałą $C_1=C_1(L,\nu(E_1))$,
\[
\|\beta(t,x)-\beta(t,y)\|^2
+\|\sigma(t,x)-\sigma(t,y)\|_{\mathrm{HS}}^2
+\int_E\|h(t,x,z)-h(t,y,z)\|^2\dd\nu(z)
\le C_1\|x-y\|^2.
\]
Zastosowanie tych samych nierówności przy $y=0$ daje
oszacowanie wzrostu
$\|\beta(t,x)\|^2+\|\sigma(t,x)\|_{\mathrm{HS}}^2+
\int_E\|h(t,x,z)\|^2\dd\nu(z)
\le C_2(1+\|x\|^2)$.

\emph{Krok 2: iteracje.} Ustalmy $X^{(0)}_t=x$ i
zdefiniujmy rekurencyjnie $X^{(n+1)}$ prawą stroną
\eqref{eq:jump-sde-comp} z $X^{(n)}_{s-}$ w miejscu $X_{s-}$.
Indukcja pokazuje, że każda iteracja jest adaptowana i càdlàg,
a wszystkie całki są kwadratowo całkowalne. Istotnie,
ze wzrostu liniowego oraz z nierówności maksymalnych
ciągłych i skokowych martyngałów otrzymujemy
\[
\E\sup_{s\le t}\|X^{(n+1)}_s\|^2
\le C_3\left(1+\|x\|^2+
 \int_0^t\E\sup_{u\le r}\|X^{(n)}_u\|^2\dd r\right),
\]
gdzie $C_3$ zależy od $T$ i danych, ale nie od $n$.

\emph{Krok 3: oszacowanie różnic.} Niech
$D^{(n)}=X^{(n+1)}-X^{(n)}$ oraz
$F_n(t)=\E\sup_{s\le t}\|D^{(n)}_s\|^2$.
Dla $n\ge1$ stosujemy kolejno
$\|u+v+w\|^2\le3(\|u\|^2+\|v\|^2+\|w\|^2)$,
$\sup_{s\le t}\|\int_0^s a_r\dd r\|^2
\le t\int_0^t\|a_r\|^2\dd r$, nierówność Dooba
oraz izometrie całek Wienera i Poissona. Dostajemy
\[
 F_n(t)\le C_T\int_0^t F_{n-1}(r)\,\dd r,
 \qquad C_T=3TC_1+12C_1.
 \tag{13.5}\label{eq:jump-picard-est}
\]
Pierwsza różnica spełnia $F_0(t)\le A$ dla $0\le t\le T$,
z pewną skończoną $A$ wynikającą z \eqref{eq:jump-at-zero}.
Indukcja całkowa prowadzi do
\[
 F_n(t)\le A\frac{(C_Tt)^n}{n!}.
 \tag{13.6}\label{eq:jump-picard-factorial}
\]
Tutaj dokładnie pojawia się silnia: każde kolejne
całkowanie po odcinku $[0,t]$ zwiększa wykładnik $t$
i dzieli oszacowanie przez następny indeks.

\emph{Krok 4: granica i równanie całkowe.} Z oszacowania
\eqref{eq:jump-picard-factorial} otrzymujemy
\[
\sum_{n\ge0}\E\sup_{t\le T}\|D^{(n)}_t\|
\le \sqrt A\sum_{n\ge0}\frac{(C_TT)^{n/2}}{\sqrt{n!}}<\infty.
\]
Twierdzenie o zbieżności monotonicznej daje więc
$\sum_n\sup_{t\le T}\|D^{(n)}_t\|<\infty$ prawie na pewno.
Iteracje mają jednostajną granicę $X$, której trajektorie
są càdlàg; proces pozostaje adaptowany.
Jednocześnie iteracje zbiegają w normie
$\E\sup_{t\le T}\|\cdot\|^2$ (nierówność Minkowskiego
i zbieżny szereg pierwiastków $F_n(T)$).
Warunek Lipschitza i oszacowania maksymalne zapewniają
zbieżność każdej z trzech całek w \eqref{eq:jump-sde-comp}
w tej samej normie. Przechodząc do granicy, otrzymujemy
tożsamość procesów, a Lemat~\ref{lem:rewrite-jump-sde}
przywraca pierwotną postać \eqref{eq:jump-sde}.

\emph{Krok 5: jednoznaczność.} Dla dwóch rozwiązań $X,Y$
ze wspólnym szumem, których normy maksymalne mają skończony
drugi moment, stosujemy \eqref{eq:jump-picard-est}
bez zmiany argumentu, zastępując $D^{(n)}$ przez $X-Y$.
W przypadku dowolnych rozwiązań o lokalnie całkowalnych
charakterystykach stosujemy najpierw wspólne czasy
zatrzymania i Lemat~\ref{lem:jump-exit-audit} podany niżej,
a następnie usuwamy lokalizację. Dla rozwiązań o
skończonym momencie maksymalnym wynika stąd
$G(t):=\E\sup_{s\le t}\|X_s-Y_s\|^2
\le C_T\int_0^tG(r)\dd r$.
Lemat Gronwalla daje $G=0$. Równość w normie supremum
daje nierozróżnialność, a nie tylko równość w ustalonej
chwili.
\end{proof}

\begin{uwaga}[Co jeżeli duże skoki nie mają drugiego momentu?]
Metoda Picarda powyżej działa w $L^2$, dlatego wymaga
\eqref{eq:jump-at-zero} dla $\Gamma$. Jeżeli $\nu(E_1)<\infty$,
można zamiast tego rozwiązywać równanie między kolejnymi dużymi
skokami, a w chwili skoku aktualizować stan według
$X_s=X_{s-}+\Gamma(s,X_{s-},Z_s)$. W tej konstrukcji
nie potrzeba momentu drugiego znacznika, ale trzeba
sprawdzić nieeksplozję i jednoznaczność dynamiki między
skokami. Nie jest to już wniosek z podanej izometrii $L^2$.
\end{uwaga}

\section{Zależność od danych, lokalizacja i nieeksplozja}

\begin{twierdzenie}[Stabilność względem stanu początkowego]
\label{th:jump-stability}
Przy założeniach Twierdzenia~\ref{th:jump-picard} niech
$X^x,X^y$ będą rozwiązaniami napędzanymi tymi samymi
$W,N$, z warunkami początkowymi $x,y\in\R^d$.
Istnieje stała $C_T$ zależna od współczynników i $T$, taka że
\[
 \E\sup_{t\le T}\|X_t^x-X_t^y\|^2\le C_T\|x-y\|^2.
 \tag{13.7}\label{eq:jump-stability}
\]
\end{twierdzenie}
\begin{proof}
Odejmując oba równania, otrzymujemy początkowy wyraz
$x-y$ oraz całki z różnic współczynników. Oszacowania
z dowodu Twierdzenia~\ref{th:jump-picard} dają
$G(t)\le4\|x-y\|^2+C\int_0^tG(s)\dd s$.
Lemat Gronwalla dowodzi $G(T)\le4e^{CT}\|x-y\|^2$.
\end{proof}

\begin{twierdzenie}[Warunek lokalny i brak eksplozji]
\label{th:jump-local-explosion}
Niech współczynniki równania \eqref{eq:jump-sde-comp}
będą lokalnie Lipschitzowskie w następującym sensie:
dla każdego $R>0$ istnieje $L_R<\infty$, dla którego
odpowiednik \eqref{eq:jump-lip} zachodzi przy
$\|x\|,\|y\|\le R$. Jeżeli dodatkowo
\[
 \|\beta(t,x)\|^2+\|\sigma(t,x)\|_{\mathrm{HS}}^2
 +\int_E\|h(t,x,z)\|^2\dd\nu(z)
 \le K(1+\|x\|^2)
 \quad(t\in[0,T],\ x\in\R^d),
 \tag{13.8}\label{eq:jump-growth}
\]
to istnieje jedyne globalne silne rozwiązanie równania
\eqref{eq:jump-sde-comp} na $[0,T]$, a
$\E\sup_{t\le T}\|X_t\|^2\le C_{K,T}(1+\|x\|^2)$.
\end{twierdzenie}
\begin{proof}
Niech $\pi_R$ będzie rzutem na domkniętą kulę promienia
$R$. Współczynniki $\beta(t,\pi_Rx)$,
$\sigma(t,\pi_Rx)$, $h(t,\pi_Rx,z)$ są globalnie
Lipschitzowskie, ponieważ rzut jest $1$-Lipschitzowski,
a argumenty pozostają w kuli promienia $R$.
Twierdzenie~\ref{th:jump-picard} daje rozwiązanie $X^R$.
Dwa rozwiązania z różnymi promieniami pokrywają się
(do chwili wyjścia z mniejszej kuli) przez jednoznaczność
z zatrzymaniem. Przy skoku równość obowiązuje także
w chwili wyjścia, gdyż współczynniki skokowe oblicza się
w stanie \emph{przed} skokiem.

Z nierówności maksymalnych i \eqref{eq:jump-growth},
niezależnie od $R$, otrzymujemy
\[
 \E\sup_{t\le T}\|X^R_t\|^2
 \le C(1+\|x\|^2)+C\int_0^T
 \E\sup_{u\le s}\|X^R_u\|^2\dd s
 \le C_{K,T}(1+\|x\|^2).
\]
Dla $\tau_R=\inf\{t\le T:\|X^R_t\|\ge R\}\wedge T$
nierówność Markowa daje
$\Prob(\tau_R<T)\le C_{K,T}(1+\|x\|^2)/R^2$.
Lewa strona dąży do zera, zatem ze zgodności lokalnych
rozwiązań otrzymujemy jedno rozwiązanie na całym $[0,T]$.
Wynik o momencie przechodzi do granicy przez lemat Fatou.
\end{proof}

\begin{uwaga}[Dlaczego lokalność nie wystarcza sama w sobie?]
Dla deterministycznego równania $x'(t)=x(t)^2$, $x(0)=1$,
wyraz szumowy jest równy zero, a dryf jest lokalnie
Lipschitzowski. Rozwiązanie $x(t)=(1-t)^{-1}$ eksploduje
w chwili $1$. Jest to szczególny przypadek SDE, więc
nie wolno wyciągać globalnego istnienia wyłącznie z
lokalnej regularności współczynników.
\end{uwaga}


\section{Lokalizacja i skok opuszczający obszar: uzupełnienie dowodowe}

\begin{lemat}[Zgodność rozwiązań także w chwili skoku wyjściowego]
\label{lem:jump-exit-audit}
Niech $X,Y$ będą dwoma adaptowanymi rozwiązaniami càdlàg
równania \eqref{eq:jump-sde-comp}, napędzanymi tą samą parą
$(W,N)$, z tym samym stanem początkowym. Załóżmy, że
współczynniki spełniają warunek Lipschitza na każdej kuli
w normie $\|\cdot\|^2+\|\cdot\|_{\mathrm{HS}}^2+L^2(\nu)$.
Dla $R>\|x\|$ oznaczmy
$\tau_R=\inf\{t\ge0:\|X_t\|\vee\|Y_t\|\ge R\}\wedge T$.
Wówczas $X_{t\wedge\tau_R}=Y_{t\wedge\tau_R}$ dla
wszystkich $t\in[0,T]$ prawie na pewno.
\end{lemat}
\begin{proof}
Ponieważ $X,Y$ są càdlàg, dla $s\le\tau_R$ zachodzi
$\|X_{s-}\|,\|Y_{s-}\|\le R$, również gdy w chwili
$s=\tau_R$ następuje skok przekraczający promień $R$.
Zatrzymujemy różnicę równań w chwili $\tau_R$. W całce
Poissona po $(0,t\wedge\tau_R]$ uwzględniony jest skok
w chwili $\tau_R$; jego współczynnik jest jednak obliczany
w stanie lewostronnym i nadal spełnia lokalne oszacowanie
Lipschitza. Stosując nierówność Cauchy'ego--Schwarza,
nierówność maksymalną Dooba i izometrie obu całek,
otrzymujemy po standardowej lokalizacji całkowalności
\[
 G_R(t):=\E\sup_{u\le t}
 \|X_{u\wedge\tau_R}-Y_{u\wedge\tau_R}\|^2
 \le C_{R,T}\int_0^t G_R(s)\,\dd s.
\]
Jeżeli zatrzymane całki nie mają z góry skończonego
momentu, najpierw zatrzymujemy je po przekroczeniu poziomu
$n$ przez ich kwadratowe charakterystyki, stosujemy to samo
oszacowanie i dopuszczamy $n\to\infty$; na odcinku do
$\tau_R$ charakterystyki są skończone prawie na pewno.
Lemat Gronwalla daje $G_R=0$. W ten sposób nie gubimy
skoku w chwili wyjścia, co byłoby błędem przy
naśladowaniu argumentu dla trajektorii ciągłych.
\end{proof}

\begin{uwaga}[Które założenie daje globalność?]
Sama lokalna lipschitzowskość zapewnia rozwiązanie do chwili
eksplozji. Oszacowanie liniowego wzrostu
\eqref{eq:jump-growth} daje ograniczenie momentu niezależne
od promienia obcięcia, a zatem $\Prob(\tau_R<T)\to0$.
Oba etapy są logicznie odrębne: zgodność rozwiązań
zatrzymanych nie oznacza jeszcze, że czasy wyjścia
zbiegają do nieskończoności.
\end{uwaga}

\section{Zadania rozwiązane i przykłady interpretacyjne}

\begin{zadanie}[Skokowy proces Ornsteina--Uhlenbecka ze znacznikami]
Niech $a,\lambda>0$, $N$ będzie miarą Poissona na
$(0,\infty)\times\R$ o intensywności
$\lambda\,\dd t\,\dd\rho(z)$, gdzie $\rho$ jest
rozkładem prawdopodobieństwa o skończonym drugim momencie.
Rozwiązać
\[
 X_t=x-a\int_0^tX_s\,\dd s
 +\int_{(0,t]\times\R}z\,\widetilde N(\dd s,\dd z),
\]
a następnie wyznaczyć średnią, wariancję i granicę wariancji.

\emph{Rozwiązanie.} Mnożymy równanie przez deterministyczny
czynnik całkujący $e^{at}$. Nie ma dodatkowego członu
kowariacji, gdyż $e^{at}$ ma skończoną wariację i jest ciągły.
Stąd
\[
 X_t=xe^{-at}+\int_{(0,t]\times\R}
 e^{-a(t-s)}z\,\widetilde N(\dd s,\dd z).
\]
Z kompensacji wynika $\E X_t=xe^{-at}$, natomiast z
izometrii Poissona
\[
 \Var(X_t)=\lambda\E_\rho Z^2\int_0^te^{-2a(t-s)}\dd s
 =\frac{\lambda\E_\rho Z^2}{2a}(1-e^{-2at}).
\]
Wariancja zbiega do $\lambda\E_\rho Z^2/(2a)$, ale
nie oznacza to automatycznie gaussowskości rozkładu
stacjonarnego.
\end{zadanie}

\begin{zadanie}[Intensywność zależna od stanu przez rozrzedzanie]
Niech $\overline N$ będzie miarą Poissona na
$(0,\infty)\times(0,\Lambda]$ o intensywności
$\dd t\,\dd u$, gdzie $0<\Lambda<\infty$, a
$\lambda:\N_0\to[0,\Lambda]$ dowolną funkcją.
Zdefiniujmy proces zliczający wzorem
\[
 X_t=x+\int_{(0,t]\times(0,\Lambda]}
  \ind{\{u\le\lambda(X_{s-})\}}\,\overline N(\dd s,\dd u),
 \qquad x\in\N_0.
\]
Wyjaśnić jego konstrukcję i obliczyć prawdopodobieństwo
jednego przejścia na krótkim przedziale.

\emph{Rozwiązanie.} Na $[0,T]$ miara $\overline N$ ma
skończenie wiele punktów prawie na pewno. Porządkujemy
je chronologicznie; w chwili punktu $(s,u)$ wykonujemy skok
$X_s=X_{s-}+1$, gdy $u\le\lambda(X_{s-})$, a w przeciwnym
razie nie zmieniamy stanu. W ten sposób otrzymujemy
jedyny proces skokowy, niekorzystający z informacji przyszłej.
Na krótkim przedziale $[t,t+h]$, pod warunkiem $X_t=k$,
prawdopodobieństwo zaakceptowania dokładnie jednego punktu
wynosi $\lambda(k)h+O(h^2)$, a prawdopodobieństwo co najmniej
dwóch akceptacji wynosi $O(h^2)$, jednostajnie w $k$.
Istotnie, liczba wszystkich punktów kandydackich ma
rozkład Poissona z parametrem $\Lambda h$, więc zdarzenie
dwóch lub więcej ma prawdopodobieństwo $O(h^2)$; przy
jednym kandydacie znacznik jest jednostajny i daje
prawdopodobieństwo akceptacji $\lambda(k)/\Lambda$.
Generator to $Lf(k)=\lambda(k)(f(k+1)-f(k))$.
\end{zadanie}

\begin{przyklad}[Dlaczego rozwiązania nie muszą zachowywać porządku?]
W całce skokowej sam warunek Lipschitza nie oznacza
porównywania trajektorii. Załóżmy, że w równaniu
$\dd X_t=-2X_{t-}\,\dd N_t$ mamy dwa warunki
początkowe $x=0$ i $y=1$. Przed pierwszym skokiem
$X^0=0$, $X^1=1$, lecz po nim $X^0=0$, $X^1=-1$.
Porządek został odwrócony. Aby uzyskać zasadę porównania,
trzeba dodatkowo kontrolować monotoniczność przekształceń
$x\mapsto x+\Gamma(t,x,z)$ wywoływanych przez skoki.
\end{przyklad}

\section{Związek z generatorami i własnością Markowa}

\begin{stwierdzenie}[Własność Markowa dla jednorodnych współczynników]
\label{prop:jump-markov}
Przy hipotezach Twierdzenia~\ref{th:jump-picard}, jeśli
$b,\sigma,\gamma,\Gamma$ nie zależą jawnie od czasu,
rodzina rozwiązań $(X^x)_{x\in\R^d}$ wyznacza
jednorodny proces Markowa o półgrupie
$P_tf(x)=\E[f(X_t^x)]$ dla każdej ograniczonej funkcji
borelowskiej $f$. Co więcej, proces ma mocną własność Markowa.
\end{stwierdzenie}
\begin{proof}
Iteracje Picarda są mierzalnymi funkcjonałami stanu
początkowego oraz przyrostów pary $(W,N)$; ich granica
jest mierzalna, a jednoznaczność ścieżkowa zapewnia
zgodność rozwiązań przesuniętych w czasie. Przy ustalonej
chwili $s$ przyszłe przyrosty pary $(W,N)$ są niezależne
od $\calF_s$. Stąd dla ograniczonej $f$
$\E[f(X_{s+t}^x)\mid\calF_s]=P_tf(X_s^x)$.
Dla ograniczonego czasu zatrzymania $\tau$ zastępujemy
$s$ przez przybliżenia $\tau_n$ o skończenie wielu
wartościach, z $\tau_n\downarrow\tau$. Prawostronna
ciągłość trajektorii, stabilność w stanie początkowym
i ciągłość stochastyczna w czasie pozwalają przejść
do granicy dla ograniczonych ciągłych $f$.
Dla dowolnej ograniczonej borelowskiej $f$ uogólniamy
wynik przez twierdzenie o klasie monotonicznej.
\end{proof}

\begin{uwaga}[Od jakiej informacji zależy przyszłość?]
Mocna własność Markowa oznacza, że po czasie zatrzymania
$\tau$ dalsza ewolucja zależy od aktualnego stanu
$X_\tau$, a nie od całej przeszłej trajektorii. Nie znaczy
to, że wartości $X_t$ w różnych chwilach są niezależne.
Właśnie zależność od $X_{s-}$ pozwala modelować
sprzężenie bieżącej dynamiki z losowymi rozmiarami skoków.
\end{uwaga}

\section{Podsumowanie i przygotowanie do rozwiązań słabych}
Metoda Picarda wraz z izometrią całki skompensowanej
zapewnia silne rozwiązania skokowych SDE przy globalnym
warunku Lipschitza. Lokalizacja poszerza zakres do
współczynników lokalnie Lipschitzowskich o liniowym
wzroście. Te wyniki nie odpowiadają jednak na pytanie,
czy można skonstruować rozwiązanie, \emph{wybierając również}
przestrzeń probabilistyczną i proces napędzający, ani czy
jednoznaczność rozkładu pociąga za sobą jednoznaczność
trajektorii. Temu poświęcimy rozdział następny.

% ===========================================================================
% ROZDZIAŁ 14 -- Słabe rozwiązania, jednoznaczność co do prawa, Yamada--Watanabe
% ===========================================================================
